\documentclass[11pt,a4paper]{article}
\usepackage[T1]{fontenc}
\usepackage[utf8]{inputenc}
\usepackage[french]{babel}
\usepackage{lmodern}
\usepackage{amsmath,amssymb}
\usepackage[margin=2.6cm]{geometry}
\usepackage{booktabs}
\usepackage{enumitem}
\usepackage{hyperref}
\usepackage{xcolor}
\hypersetup{colorlinks=true,linkcolor=blue!50!black,urlcolor=blue!50!black}

\newcommand{\NW}{\mathrm{NW}}
\newcommand{\E}{\mathbb{E}}
\newcommand{\demontre}[1]{\textbf{[démontré]}~#1}
\newcommand{\plausible}[1]{\textbf{[plausible]}~#1}
\newcommand{\hypothese}[1]{\textbf{[hypothèse]}~#1}
\newcommand{\arbitraire}[1]{\textbf{[choix arbitraire]}~#1}
\newcommand{\artefact}[1]{\textbf{[artefact possible]}~#1}

\title{Lecture critique du document de conception M2\\
\large « Conception d'un modèle à distribution Boltzmann--Pareto émergente
par auto-organisation critique »}
\author{Agent de réplication (workspace \texttt{m2\_fable})}
\date{3 juillet 2026}

\begin{document}
\maketitle

\begin{abstract}
Ce rapport analyse le document de conception du modèle M2 comme une
\emph{proposition de recherche à tester}, non comme un acquis. Nous en donnons
un résumé fidèle (\S\ref{sec:resume}), une reconstruction mathématique
(\S\ref{sec:math}), une liste des hypothèses fortes avec leur statut
épistémique (\S\ref{sec:hypotheses}) et une critique scientifique
(\S\ref{sec:critique}). Conclusions principales : (i)~le modèle est
remarquablement discipliné dans sa traque des paramètres cachés, mais trois
conventions non testées (moyenne géométrique des taux, appariement
max/min, transferts au prorata) et la fenêtre de naissance
$(w_0, d_0)$ constituent des degrés de liberté réels ; (ii)~les résultats
[E7c] présentés comme « pile complète M2 » proviennent d'un prototype qui
\emph{dévie de la spécification M2} sur au moins deux règles (volume de prêt,
sort des créances de la faillie) — la réplication n'est donc pas garantie ;
(iii)~le corps exponentiel est explicitement non acquis et la fenêtre de
naissance ressemble à un réglage fin déguisé, ce que le protocole §6.5 du
rapport a le mérite d'exposer à la réfutation ; (iv)~le garde-fou anti-cohorte
du rapport est nécessaire mais incomplet : nous proposons des tests
supplémentaires.
\end{abstract}

\tableofcontents

%=====================================================================
\section{Position épistémique et sources}

Le document de conception (noté ci-après [CONC]) est daté du 2 juillet 2026 et
s'appuie sur : Bouchaud--Mézard 2000 [BM], Yakovenko--Rosser 2009 [YR], le
modèle historique [WIP] (\texttt{modele-27-04-WIP}), une analyse de
distributions antérieure [DISTRIB], et une chaîne d'explorations numériques
E0--E7 exécutées sur un prototype \texttt{proto\_m1.py}.

Statuts que nous distinguons systématiquement :
\demontre{} établi par une ablation reproductible ou un calcul exact ;
\plausible{} argumenté analytiquement mais non vérifié sur M2 ;
\hypothese{} de travail, à tester ;
\arbitraire{} convention de modélisation parmi d'autres possibles ;
\artefact{} numérique ou de protocole possible.

Un point de méthode s'impose d'emblée : \emph{les ablations E0--E7 portent sur
le prototype M1, pas sur M2}. Le prototype conserve deux règles héritées du
[WIP] que M2 supprime (volume de prêt $q = \min(w_l/2,\ \theta\, q_{\max})$
avec $\theta = 0{,}35$, au lieu de la règle symétrique
$q = \min(q_{\mathrm{dem}}, q_{\mathrm{off}})$), et sa règle de faillite
\emph{annule} les contrats où la faillie est prêteuse au lieu de les
transférer au prorata (fonction \texttt{\_fail} du prototype, commentaire
« gain pour lui » — qui décrit d'ailleurs l'annulation des dettes des
\emph{emprunteurs} de la faillie, un troisième écart de comptabilité).
Les verdicts qualitatifs des ablations (nécessité de l'asymétrie
nominal/réel, du coût $r+\delta$, de $d_0$) restent informatifs ; les valeurs
numériques de [E7c] (queue $\alpha \approx 2{,}8$ stable, $N^* \approx 1760$)
ne sont pas des prédictions pour M2 au sens strict — le rapport le dit
lui-même (§8, incertitude 5), mais l'affirmation « trois des quatre cibles
atteintes » du résumé exécutif doit être lue avec cette réserve.

%=====================================================================
\section{Résumé fidèle du modèle M2}
\label{sec:resume}

\subsection{État d'une entité}
Une entité vivante $i$ porte :
\begin{itemize}[nosep]
\item un capital réel scalaire $w_i \ge 0$ (une seule colonne de bilan) ;
\item une dette d'amorçage nominale $d_0$, constante universelle de naissance,
  jamais remboursée, jamais dépréciée, détenue par personne ;
\item des contrats de prêt la liant à d'autres entités.
\end{itemize}

\subsection{Contrats de prêt}
Un contrat est un quadruplet (prêteur $\ell$, emprunteur $b$, principal
nominal $q > 0$, taux $r$). Il est \emph{perpétuel} : jamais amorti, $q$
constant de la création à l'annulation. Le principal est nominal : il ne se
déprécie pas, contrairement au capital réel. Dérivés : créances
$C_i = \sum_{\ell = i} q$, dettes $D_i = \sum_{b = i} q$.

\subsection{Valeur nette et revenu}
\begin{equation}
\NW_i = w_i + C_i - D_i - d_0
\end{equation}
est la variable cible de la forme Boltzmann--Pareto. Le revenu est
$y_i = \alpha\sqrt{w_i} + (\text{intérêts reçus})_i$, le revenu net
$y_i - (\text{intérêts payés})_i$ ; capital brut et revenu sont mesurés à
titre diagnostique.

\subsection{Séquence d'un pas de temps (ordre fixé, 7 phases)}
\begin{enumerate}[nosep]
\item \textbf{Naissances} : $n \sim \mathrm{Poisson}(\lambda)$ ; dotation
  $w_0$, dette $d_0$, donc $\NW$ initial $\varepsilon_0 = w_0 - d_0 > 0$
  petit. $N(0) = 0$ (pas de cohorte fondatrice).
\item \textbf{Choc multiplicatif commun en loi} :
  $w_i \leftarrow w_i e^{\eta_i}$, $\eta_i \sim
  \mathcal{N}(-\sigma^2/2, \sigma^2)$ i.i.d., donc $\E[e^\eta] = 1$. Frappe le
  capital réel, pas les contrats.
\item \textbf{Extraction} : $w_i \leftarrow w_i + \alpha\sqrt{w_i}$,
  $\alpha \equiv 1$.
\item \textbf{Service des intérêts} : chaque contrat, l'emprunteur paie $rq$
  au prêteur (transfert conservatif) ; paiement partiel possible (tout le $w$
  disponible), alors \emph{défaut} $\to$ faillite en phase 7.
\item \textbf{Dépréciation} : $w_i \leftarrow (1-\delta) w_i$ ; contrats et
  $d_0$ non dépréciés (asymétrie constitutive).
\item \textbf{Marché du crédit} : $\lfloor N/k \rfloor$ rounds ; par round,
  tirage uniforme sans remise de $k$ vivantes ; prêteur $=$ $w$ maximal,
  emprunteur $=$ $w$ minimal ; taux $r = \sqrt{r^*_\ell\, r^*_b}$ avec
  $r^* = \alpha/(2\sqrt{w})$ ; coût effectif $\rho = r + \delta$ ; cible
  commune $w^*(\rho) = (\alpha/2\rho)^2$ ; demande
  $q_{\mathrm{dem}} = \max(0, w^* - w_b)$, offre
  $q_{\mathrm{off}} = \max(0, w_\ell - w^*)$, volume
  $q = \min(q_{\mathrm{dem}}, q_{\mathrm{off}})$ ; si $q>0$, transfert
  conservatif du principal et création du contrat.
\item \textbf{Faillites en cascade} : critère $\NW_i < 0$ ou défaut en
  phase 4. À la faillite : contrats empruntés annulés (perte sèche des
  prêteurs) ; contrats prêtés \emph{transférés aux créanciers au prorata} ;
  capital résiduel réparti aux créanciers au prorata ; $d_0$ éteinte. Itérer
  jusqu'à stabilité. La mort est la seule sortie.
\end{enumerate}

\subsection{Fenêtre de naissance}
Contrainte de cohérence quantitative établie par [E6]--[E7] :
\begin{equation}
\label{eq:fenetre}
\frac{1}{(\sigma+\delta)^2} \;\lesssim\; w_0 \;<\; w^*(\bar r + \delta),
\qquad \varepsilon_0 \lesssim \sigma w_0 .
\end{equation}
Borne basse : la diffusion doit dominer le drift $\sqrt{w}$ à la naissance
(sinon aucun entrant ne meurt, [E7a]) ; borne haute : l'entrante doit naître
demandeuse de crédit (sinon le marché meurt, [E6]) ; marge $\varepsilon_0$ de
l'ordre d'un choc (sinon mortalité infantile négligeable, [E7b]).

\subsection{Population ouverte}
Bornage endogène attendu : $N^* = \lambda\,\E[\text{vie}]$, chaque position
(entrante, endettée, créancière) ayant une espérance de vie finie. Symptôme
d'échec : $N(t)$ affine $=$ sous-population immortelle.

\subsection{Paramètres irréductibles}
$\alpha = 1$ (unité de flux), $\delta$ (unité de temps), $\lambda$ (taille
$N^*$), $k$ (connectivité, contrôle SOC), $\sigma$ (variance du choc commun),
$(w_0, d_0)$ ou $(w_0, \varepsilon_0)$ (bilan de naissance). Critère
d'acceptation « sans réglage fin » : la \emph{forme} Boltzmann--Pareto doit
être robuste sur $(\sigma, \varepsilon_0/w_0, k)$, seuls $(T, \mu, X_0)$
pouvant bouger de façon lisse.

%=====================================================================
\section{Reconstruction mathématique}
\label{sec:math}

\subsection{Dynamique individuelle}
Entre événements de crédit, le capital suit la version discrétisée de
\begin{equation}
\frac{dw}{dt} = -\delta w + \alpha\sqrt{w} + c_{\mathrm{eff}},
\qquad c_{\mathrm{eff}} = \text{intérêts reçus} - \text{payés}
\pm \text{principal} \pm \text{pertes},
\end{equation}
plus le choc de la phase 2 : $\ln w$ reçoit $\eta - \sigma^2/2$ par pas.
En autarcie ($c_{\mathrm{eff}} = 0$, $\sigma = 0$), point fixe commun
$x_+ = (\alpha/\delta)^2$ (pour $\alpha=1$, $\delta = 0{,}05$ : $x_+ = 400$).
Ce point fixe est « l'ennemi » : tout le modèle sert à tenir la population
loin de lui [E0].

Un calcul utile pour la suite : la composition exacte d'un pas sans crédit
donne
\begin{equation}
\label{eq:map}
w_{t+1} = (1-\delta)\left( w_t e^{\eta_t} + \alpha\sqrt{w_t e^{\eta_t}} \right),
\end{equation}
dont le point fixe déterministe ($\eta = 0$) est
$\hat w = \alpha^2 (1-\delta)^2 / \delta^2$
(résolution de $\hat w = (1-\delta)(\hat w + \sqrt{\hat w})$, soit
$\sqrt{\hat w}\,\delta = (1-\delta)$, $\hat w = ((1-\delta)/\delta)^2 = 361$
pour $\delta = 0{,}05$) — légèrement sous le $x_+ = 400$ du temps continu.
L'écart n'a pas d'importance qualitative mais rappelle que \emph{les repères
quantitatifs du rapport sont des repères de temps continu appliqués à une
carte discrète} ; $w^*(\rho)$ relève de la même approximation.

\subsection{Condition marginale et cible de capital}
Le rendement marginal brut de l'extraction est
$\partial_w(\alpha\sqrt{w}) = \alpha/(2\sqrt{w}) =: r^*(w)$, décroissant.
Une unité de capital empruntée coûte $r$ d'intérêt \emph{et} $\delta$ de
fonte (la dette, nominale, ne fond pas) : coût effectif $\rho = r+\delta$.
L'égalisation $r^*(w) = \rho$ donne la cible commune
\begin{equation}
w^*(\rho) = \left(\frac{\alpha}{2\rho}\right)^2 .
\end{equation}
Pour $\bar r \approx 0{,}05$--$0{,}07$ et $\delta = 0{,}05$,
$w^* \approx 17$--$25$ ; avec la convention du rapport
($\bar r + \delta \approx 0{,}075$), $w^* \approx 44$. La règle de marché
« viser $w^*$ » est l'optimalité myope à un pas, correctement écrite.
Remarque critique : $\rho$ ignore le risque de défaut du côté du prêteur et
le risque de faillite du côté de l'emprunteur — c'est une optimalité
\emph{sans anticipation}, choix de modélisation cohérent avec le cadre
(agents à règles fixes) mais qu'il faut nommer : le taux ne contient aucune
prime de risque, alors que le défaut est le mécanisme central du modèle.

\subsection{Fokker--Planck sur la valeur nette}
Le squelette analytique est [YR] éq.~(23--25) : pour
$A(X) = A_0 + aX$, $B(X) = B_0 + bX^2$,
\begin{equation}
P(X) \propto \frac{1}{B(X)} \exp\left( -\int^X \frac{A(X')}{B(X')} dX' \right)
\;\Longrightarrow\;
\begin{cases}
P(X) \propto e^{-X/T}, & X \ll X_0, \quad T = B_0/A_0,\\[2pt]
P(X) \propto X^{-1-\mu}, & X \gg X_0, \quad \mu = 1 + a/b,
\end{cases}
\end{equation}
avec croisement $X_0 = \sqrt{B_0/b}$. L'argument du rapport : près du
plancher, les flux qui bougent $\NW$ (marge $\varepsilon_0$, intérêts $rq$
dimensionnés par le marché, principal, extraction nette évaluée près de
$w^*$) sont bornés et indépendants de $X$ (régime additif) ; pour une
créancière établie, chocs $\sigma w$, revenus $\bar r C$ et pertes de cascade
sont $\propto X$ (régime multiplicatif).

\textbf{Point faible de l'argument} (repris en \S\ref{sec:critique}) :
l'hypothèse « $A_0 > 0$ » (drift net \emph{vers le bas} dans le corps) n'est
pas démontrée ; et surtout l'extraction $\alpha\sqrt{w} - \delta w$ n'est pas
un bruit mais un \emph{rappel déterministe} vers $w^*$ ou $x_+$ : un
processus d'Ornstein--Uhlenbeck avec rappel donne un corps gaussien autour de
la cible, pas une exponentielle. Le corps exponentiel exige que le bruit
additif de crédit domine le rappel — le rapport le reconnaît (incertitude 1)
mais ne quantifie pas le rapport des deux échelles. C'est, selon nous, la
raison structurelle pour laquelle [E7c] trouve un corps gamma/log-normal
(unimodal à mode positif) plutôt qu'exponentiel.

\subsection{Queue de Kesten}
Pour une établie : $X_{t+1} = \Lambda_t X_t + \beta_t$ avec
\begin{equation}
\Lambda_t = 1 + \eta_t - \delta h_t + \bar r c_t - \ell_t,
\end{equation}
$h = w/X$, $c = C/X$, $\ell$ pertes de défaut relatives. Queue Pareto dès que
$\E[\ln\Lambda] < 0 < P(\Lambda > 1)$, exposant racine de
$\E[\Lambda^\mu] = 1$ [BM éq.~8--9]. Repère sans crédit ($h=1$, $c=\ell=0$),
$\eta$ log-normal : $\mu \approx 1 + 2\delta/\sigma^2$ (pour
$\delta = 0{,}05$, $\sigma = 0{,}25$ : $\mu = 2{,}6$, exposant de pdf
$3{,}6$). Le crédit \emph{abaisse} $\mu$ ($\bar r c$ compense $\delta h$), les
cascades le \emph{relèvent} ($\ell$). Vérification de cohérence interne : la
dérivation $\mu = 1 + 2\delta/\sigma^2$ vient de
$\E[\Lambda^\mu] = 1$ avec $\ln\Lambda \approx \eta - \delta$,
$\E[e^{\mu\eta}] = e^{\mu(\mu-1)\sigma^2/2}$, d'où
$\mu(\mu-1)\sigma^2/2 = \mu\delta$ soit $\mu = 1 + 2\delta/\sigma^2$ —
correcte.

\textbf{Réserve} : pour une créancière, le choc $\eta$ ne frappe que la part
réelle $h$ de $X$, donc la variance multiplicative effective est
$\sigma^2 h^2 < \sigma^2$ ; le canal cascade $\ell$ est intermittent et
corrélé entre grosses créancières (choc systémique, pas i.i.d.) — le cadre
Kesten i.i.d. est une approximation dont la validité doit être vérifiée
empiriquement (stationnarité de $\mu$ par fenêtres, protocole §6.2).

\subsection{L'hypothèse auto-critique H2}
La boucle proposée : queue plus lourde $\Rightarrow$ crédit plus concentré
$\Rightarrow$ cascades plus grandes $\Rightarrow$ $\ell \uparrow$
$\Rightarrow$ queue plus légère (et réciproquement). Prédictions testables :
$\mu$ stationnaire par fenêtres ; $\mu$ insensible à $\lambda, w_0$ et à $k$
dans la phase SOC ; tailles de cascade à queue lourde ; raidissement continu
de $\mu$ quand $\sigma \to 0$. \hypothese{} — le rapport est explicite
là-dessus ; aucun mécanisme quantitatif (gain de boucle, temps de relaxation)
n'est donné, et rien ne garantit que le point fixe de la rétroaction soit
dans la zone « queue lourde observable » plutôt que $\mu \gg 3$
(indiscernable d'une log-normale sur $N \sim 10^3$).

%=====================================================================
\section{Hypothèses fortes et statuts}
\label{sec:hypotheses}

\begin{enumerate}
\item \textbf{Émergence du corps exponentiel sur $\NW$ (H1).}
  \hypothese{} Le rapport lui-même classe X1 comme « LA question ouverte » ;
  [E7c] trouve gamma/log-normale devant l'exponentielle
  ($\mathrm{med}/\mathrm{mean} = 0{,}52$--$0{,}55$ contre $0{,}69$ attendu).
  Notre analyse (\S\ref{sec:math}) suggère une raison structurelle (rappel
  déterministe vers $w^*$) qui ne disparaîtra pas en peuplant mieux le corps.
  L'issue de repli du rapport (corps [BM] $\propto e^{-(\mu-1)/X} X^{-1-\mu}$)
  est une hypothèse alternative légitime à tester en parallèle.
\item \textbf{Stabilité auto-organisée de l'exposant de Pareto (H2).}
  \hypothese{} Soutenue par un seul run long, un seul seed, sur le prototype
  déviant. Trois fenêtres avec $\alpha = 2{,}80/2{,}86/2{,}88$ est un indice,
  pas une démonstration (pas de bootstrap rapporté, dérive monotone
  $+0{,}08$ qui pourrait être un transitoire lent).
\item \textbf{Rôle suffisant des faillites en cascade comme régulateur.}
  \plausible{} mais la boucle H2 n'a jamais été observée directement
  (aucune mesure de corrélation concentration $\leftrightarrow$ cascades dans
  les logs).
\item \textbf{Absence d'effet cohorte.} \plausible{} sur le prototype
  ([E7c] : $\mathrm{corr}(\text{âge}, \log w) = 0{,}11$--$0{,}15$, queue à âge
  contrôlé). Mais mesuré sur $w$, pas sur $\NW$ ; et le renouvellement du top
  décile (sonde X2) n'a pas été exécuté sur la pile complète — le log
  \texttt{renewal\_*.log} concerne la règle naïve.
\item \textbf{Robustesse sans réglage fin.} \hypothese{} La fenêtre de
  naissance (\ref{eq:fenetre}) est étroite : [E7a] (échec), [E7b] (échec),
  [E7c] (succès) documentent précisément un \emph{calage} de $(w_0, d_0)$.
  Le rapport requalifie ce calage en « condition d'existence », analogue à
  choisir une température sous le point critique ; l'argument est défendable
  \emph{si} la forme est robuste à l'intérieur de la fenêtre — exactement ce
  que la grille §6.5 doit établir. Tant qu'elle ne l'a pas fait, la
  distinction condition d'existence / réglage fin est rhétorique.
\item \textbf{Choix de $\NW$ comme variable cible.} \arbitraire{} motivé
  (c'est la seule variable avec plancher d'absorption naturel), mais $\NW$
  est aussi la seule variable \emph{non observable} empiriquement chez [YR]
  (les données de revenu de [YR] concernent $y$). Si le corps exponentiel
  n'apparaît que sur $\NW$ et pas sur $y$, la correspondance avec la
  littérature empirique s'affaiblit.
\item \textbf{Moyenne géométrique des taux.} \arbitraire{} assumé. Notons que
  $r = \sqrt{r^*_\ell r^*_b} = \alpha / (2 (w_\ell w_b)^{1/4}\,\cdot)$ dépend
  des \emph{deux} capitaux ; la moyenne arithmétique donnerait des taux plus
  hauts (AM $\ge$ GM), donc plus de revenus d'intérêts et une queue plus
  lourde — le rapport propose d'ailleurs ce levier en X1.3. Convention à
  variante planifiée, correctement signalée.
\item \textbf{Transfert prorata (créances et capital résiduel de la
  faillie).} \arbitraire{} et \emph{potentiellement structurant} : donner les
  créances et le capital de la faillie aux créancières \emph{renforce les
  grosses} — un mécanisme de condensation directe qui pourrait fabriquer la
  queue à lui seul. Le rapport le sait (X2, variante iii « destruction
  pure ») mais ne l'a pas testé. C'est notre candidat n°1 d'explication
  alternative pour la queue.
\item \textbf{Population bornée endogènement.} \plausible{} sur le prototype
  [E1, E6, E7c]. Mécanisme clair (mortalité universelle via $d_0$). Le risque
  résiduel est une sous-population quasi immortelle à grand $w$ (mortalité du
  top $> 0$ à vérifier, critère 6.4).
\item \textbf{Choc commun « en loi » i.i.d. entre entités.} \arbitraire{}
  et sémantiquement ambigu : le rapport dit « choc multiplicatif commun »
  mais spécifie $\eta_i$ \emph{i.i.d. par entité} (le prototype tire un
  vecteur). Il s'agit d'une loi commune, pas d'un choc agrégé commun. Un vrai
  choc agrégé (même $\eta$ pour toutes au pas $t$) changerait complètement la
  physique des cascades. La spécification doit figer cette lecture.
\end{enumerate}

%=====================================================================
\section{Critique scientifique}
\label{sec:critique}

\subsection{Le modèle est-il vraiment minimal ?}
Presque. La traque des paramètres du [WIP] (18 $\to$ 6) est rigoureuse et la
plupart des suppressions sont justifiées par ablation. Trois réserves :
(i)~la \emph{fenêtre de naissance} fait de $(w_0, d_0)$ un couple contraint à
deux inégalités près — fonctionnellement, c'est un paramètre de forme calé,
même si le rapport l'assume ; (ii)~$\lfloor N/k \rfloor$ rounds par pas est
présenté comme « définition du pas de marché », mais l'intensité de churn
(rounds $\times$ taille d'échantillon / $N$) est un paramètre caché que le
rapport lui-même propose de faire varier en X1.3 — donc il en est un ;
(iii)~la règle max/min dans l'échantillon de taille $k$ couple la sélection
des contreparties à $k$ : l'écart inter-quantiles des extrêmes croît avec
$k$, donc $k$ contrôle \emph{à la fois} la connectivité et l'intensité des
gradients exploités — deux rôles pour un paramètre, ce qui compliquera
l'interprétation du balayage en $k$.

\subsection{Quels paramètres restent cachés ?}
\begin{itemize}[nosep]
\item la garde numérique $\epsilon$ (tolérances de faillite, de volume
  minimal de prêt) — inoffensive d'après [ÉLAGAGE], à vérifier sur M2 ;
\item l'\emph{ordre de service des contrats} en phase 4 quand $w$ ne couvre
  pas tous les intérêts dus : servir dans l'ordre de création, en ordre
  aléatoire ou au prorata change qui subit le défaut ; non spécifié ;
\item la \emph{participation au marché} des entités déjà en défaut au même
  pas (le prototype les exclut ; le rapport ne le dit pas) ;
\item le seuil de viabilité d'un contrat ($q$ minimal, $r$ minimal quand
  $w_\ell \approx w_b$) ;
\item $T_{\mathrm{transitoire}}$ (500 pas) et la cadence de snapshot (50) —
  paramètres de mesure, à justifier par les autocorrélations.
\end{itemize}

\subsection{Règles contestables économiquement ou mathématiquement}
(i)~Le prêteur prête \emph{tout} son excédent au-dessus de $w^*$ à une seule
contrepartie par round : aucune diversification, concentration maximale du
risque — favorable aux cascades (voulu ?) mais économiquement extrême.
(ii)~Le taux sans prime de risque face à un défaut endémique : les
créancières ne « prix-ent » jamais le risque, ce qui entretient
mécaniquement le sur-prêt — c'est peut-être le vrai moteur de la SOC, à
nommer. (iii)~$d_0$ éteinte à la faillite sans créancier : de la monnaie
nominale disparaît sans contrepartie ; la « masse de valeur nette » agrégée
$\sum \NW = W_{\mathrm{tot}} - N d_0$ n'est conservée par aucune règle —
acceptable (l'économie est ouverte) mais à surveiller dans les bilans.
(iv)~Le transfert des contrats prêtés de la faillie à ses créanciers crée
des liens entre entités qui ne se sont jamais « rencontrées », allongeant
les chaînes d'intermédiation au fil des faillites — mécanisme intéressant
mais dont la stationnarité n'est pas évidente (la profondeur des chaînes
peut-elle dériver ?).

\subsection{Mécanismes pouvant créer artificiellement une queue de Pareto}
\begin{enumerate}[nosep]
\item \emph{Condensation par prorata} : capital résiduel et créances des
  faillies vont aux plus grosses créancières (cf. \S\ref{sec:hypotheses}.8).
\item \emph{Effet cohorte résiduel} : l'ère d'amorçage ($t \lesssim 100$)
  reste anormalement favorable ; contrôlé par 6.3 mais à re-vérifier sur
  $\NW$.
\item \emph{Pooling de fenêtres} : agréger des instantanés d'un régime non
  stationnaire fabrique des queues (mixture de log-normales) ; le protocole
  s'en défend (fenêtres séparées) mais l'estimateur $x_{\min}$ qui saute
  [E1] montre que le danger est réel.
\item \emph{Choc log-normal lui-même} : avec $\sigma = 0{,}25$, un GBM à
  rappel faible produit des queues log-normales lourdes difficiles à
  distinguer d'une Pareto sur 2--3 décades ; le LR de Vuong est le bon
  outil, à condition d'être appliqué par fenêtre et par seed.
\end{enumerate}

\subsection{Mécanismes pouvant empêcher le corps exponentiel}
Le rappel déterministe $\alpha\sqrt{w} - \delta w$ vers $w^*$ (côté $w$) ;
la réinjection ponctuelle à $\varepsilon_0$ (le corps de [YR] exige une
réinjection \emph{par le bas}, mais ici toutes les naissances entrent au
même point — une masse de Dirac à $\varepsilon_0$ lissée par diffusion, pas
un flux thermique) ; le fait que les incréments de $\NW$ des petites entités
sont dominés par $rq$ et les chocs $\sigma w^*$, qui ont une échelle fixe
mais un \emph{drift} qui dépend de la position dans le cycle de vie (jeune
endettée vs établie redescendue) — le « corps » peut être une mixture de
sous-populations dynamiques même sans être une mixture d'âges.

\subsection{Les diagnostics proposés suffisent-ils ?}
Le protocole §6 est sérieux (AIC/BIC multi-familles, CSN par fenêtres,
$x_{\min}$ commun, anti-cohorte, renouvellement). Manques :
\begin{itemize}[nosep]
\item aucun test de \emph{stationnarité formelle} des séries $N(t)$,
  $W_{\mathrm{tot}}(t)$ (une pente lente peut passer pour un plateau à
  l'œil) ;
\item pas de mesure des \emph{autocorrélations temporelles} justifiant
  l'espacement de 50 pas des instantanés poolés (la vie médiane ~20--30 pas
  est un argument, pas une mesure) ;
\item le test mécanistique $T \approx B_0/A_0$ est excellent mais il faut
  préciser \emph{comment} conditionner « être dans le corps » (fenêtre de
  $\NW$ fixe, pas quantile mouvant, sinon biais de sélection) ;
\item l'alternative corps-[BM] (éq.~7) n'est pas dans l'échelle de familles
  d'ajustement — l'ajouter ;
\item aucun contrôle par \emph{modèle nul} : un GBM à rappel + plancher
  absorbant + réinjection, sans crédit, calé sur les mêmes moments — pour
  savoir ce que le crédit ajoute réellement.
\end{itemize}

\subsection{Tests supplémentaires proposés (au-delà du §6)}
(a)~variante « destruction pure » du résiduel et « annulation symétrique »
des créances de la faillie (déjà en X2, à exécuter systématiquement, pas en
recours) ; (b)~variante moyenne arithmétique des taux ; (c)~modèle nul sans
crédit ($k$ tel que le marché est vide) à $\sigma$ égal ; (d)~traçage des
trajectoires individuelles (rangs) pour mesurer la mobilité
(matrice de transition entre quintiles sur 100 pas) ; (e)~bilan comptable
par pas (conservation de $w$ total aux flux identifiés près) comme test
d'intégrité continu.

\subsection{Où le modèle risque-t-il de reproduire un artefact de cohorte ?}
Trois endroits : l'ère d'amorçage (favorable, contrôlée par exclusion du
transitoire — mais 500 pas suffisent-ils si $\E[\text{vie}]$ du top est
longue ?) ; le bas de la distribution (les naissances à $\varepsilon_0$
peuvent dominer les premiers déciles de $\NW$ — le garde-fou « établies
redescendues » de X1.1 est le bon test) ; et \emph{la queue à âge contrôlé
sur $\NW$} (jamais mesurée : [E7c] l'a fait sur $w$).

\subsection{Capital, revenu, valeur nette : distinction bien exploitée ?}
Oui dans la conception (NW = variable de Boltzmann, $y$ = variable de la
thèse énoncée, $w$ = diagnostic), mais les validations [E7c] mesurent
surtout $w$ ; l'exposant du revenu ($4{,}5$--$4{,}9$) diffère nettement de
celui de $\NW$ ($2{,}7$--$2{,}8$), ce qui contredit la prédiction §3.5
($y \approx \bar r C \propto \NW$ dans la queue $\Rightarrow$ même exposant).
Le rapport ne commente pas cet écart — c'est un signal que la queue du
revenu est portée par $\sqrt{w}$ plus longtemps que prévu, ou que la
financiarisation $c = C/X$ des grosses est plus faible qu'attendu. À
instrumenter ($h$, $c$, $\ell$ par décile).

\subsection{Modèle de richesse, de crédit ou démographique ?}
Structurellement, M2 est un \emph{modèle de crédit à démographie
absorbante} : la forme du corps dépend du flux naissance/faillite
(démographie), la queue dépend du carnet de contrats (crédit), et le
capital réel ne sert que de substrat aux chocs et à l'extraction. La
« richesse » est un épiphénomène comptable ($\NW$). Ce n'est pas un défaut —
mais cela implique que les paramètres démographiques ($\lambda$,
$\varepsilon_0$) et de crédit ($k$, règle de taux) sont les vrais leviers,
et que la robustesse doit être démontrée d'abord contre eux.

\subsection{Verdict}
Le document est une proposition de recherche de bonne qualité : falsifiable,
instrumentée, honnête sur ses échecs (corps non exponentiel dit trois fois).
Ses deux points aveugles principaux : (1)~l'écart entre le prototype ayant
produit [E7c] et la spécification M2 (volume de prêt, faillite) n'est
signalé nulle part — le risque de non-réplication est réel et c'est la
première chose que l'implémentation doit vérifier ; (2)~le mécanisme de
condensation par prorata n'est pas contrôlé alors qu'il suffit peut-être à
produire la queue. Notre implémentation traitera ces deux points comme des
expériences de première priorité.

\end{document}
